\chapter{Trajectories in exact arithmetic}

\textbf{The reproduction in this chapter is exact and the comparison to measurement is not run.} The
radii, energies and line ratios below are exact rationals, facts of the Bohr model with no digit
approximated. Whether a real atom matches them is the oracle stage, and it is named at the end of the
chapter, not carried out.

\section{The orbit is a rational}

A hydrogen-like ion is one electron on a nucleus of charge $Z$. Bohr quantizes the angular momentum
at an integer multiple of $\hbar$, and the orbit that follows has radius $n^2/Z$ in units of the Bohr
radius and energy $-Z^2/n^2$ in units of the Rydberg. Both are exact rationals, an integer over an
integer, built in
\texttt{examples/\allowbreak{}particle\_physics/\allowbreak{}1\_represent/\allowbreak{}an\_orbit\_as\_exact\_rationals.py}
with integer arithmetic done by hand, the host stand-in for the fixed-width limb integers the device
arm carries in \texttt{no\_rounding}.

The quantum number $n$ is an exact integer, so two neighboring orbits differ by
\[
\frac{(n+1)^2}{Z} - \frac{n^2}{Z} \;=\; \frac{2n+1}{Z},
\]
exactly, and no two orbits ever fall on one radius, however large $n$ grows. That is the property a
floating point orbit loses: past some $n$ the gap between two orbits falls below the last bit, and the
two large orbits merge. The exact reading keeps them apart forever, the same infinite
discrimination the shells were read with, now carrying a trajectory instead of a state.

\section{One hub, propagated exactly}

Every energy level of every hydrogen-like ion is $E(n, Z) = -\mathrm{Ry}\, Z^2/n^2$, and every radius
is $r(n, Z) = a_0\, n^2/Z$. One physical quantity sets each, the Rydberg energy and the Bohr radius,
and the whole grid of levels and radii follows from it by an exact rational identity. Compute the hub
once and the ladder is fixed to the precision of the hub, with nothing added by the arithmetic that
spreads it.

The split this draws is the split a precision program needs. The ratios $n^2/Z$ and $1/n^2$, the
integer quantum numbers, the capacities and the closures are mathematical: their uncertainty is
exactly zero and there is nothing in them to measure better. The Rydberg energy, the Bohr radius, the
fine-structure constant and the electron-to-proton mass ratio are physical: they are measured, their
uncertainty is irreducible, and they are the only floor the exact quantities stand on.

\section{The ratios carry no constant}

Within one spectral series the hub cancels. A transition from $n_2$ to $n_1$ has photon energy
proportional to $1/n_1^2 - 1/n_2^2$, and the ratio of two lines of one series is a pure rational with
no Rydberg and no meter left in it. The Balmer series, $n_1 = 2$, read as wavelength ratios to its
first line:

\begin{center}
\begin{tabular}{@{}llll@{}}
\toprule
line & upper $n$ & ratio & decimal \\
\midrule
H$\alpha$ & 3 & $1$      & $1.00000$ \\
H$\beta$  & 4 & $20/27$  & $0.74074$ \\
H$\gamma$ & 5 & $125/189$& $0.66138$ \\
H$\delta$ & 6 & $5/8$    & $0.62500$ \\
\bottomrule
\end{tabular}
\end{center}

\noindent These ratios are predictions with no measurement in them, infinitely precise, and a whole
ladder of them stands on no physical constant at all. A measured ratio meets them at the oracle, and
the meeting needs no constant either, and the comparison stays clean.

\section{What the oracle needs, and what the model is}

The absolute size of an orbit in meters, and any comparison of an absolute line to a measured
wavelength, needs the Rydberg constant and the Bohr radius, which are measured and cited, not derived.
That comparison is the oracle stage. It is not run here, because it needs the measured
configurations and the measured spectra, and those are not fetched. Naming a source is not reading
it, and no measured wavelength is quoted in this book.

The floor on a real spectral line is not the arithmetic, which stays exact, but two things above it.
One is the measured constants. The other, larger by far, is the model: the Bohr picture omits fine
structure, the Lamb shift and every later correction, so the model error dominates the constants'
uncertainty long before the arithmetic enters. The exact rationals are the model reproduced without
loss. How far the model itself stands from a measured atom is the question the oracle exists to
answer, and it is left open here.
